\section{Monitoring CCTV}
Zastosowano \CameraCount{} kamer stałych AXIS P3265-LVE w wersji obiektywu 9 mm, 1920 $\times$ 1080 px. Zakres poziomego pola widzenia 36--100$^\circ$ umożliwia przyjęte kadry. Maksymalny pobór urządzenia wynosi 10,7 W; alokacja portu PSE klasy 3 wynosi 15,4 W. Wybrano jeden model z obudową zewnętrzną także do wnętrz, aby ograniczyć liczbę części zamiennych. Jest to referencyjny model z opublikowaną kartą, nie deklaracja bieżącej dostępności sprzedażowej. PTZ nie zastosowano: wymagane kierunki mają być dostępne jednocześnie.\cite{S-CAM}

\CameraTable

Punkty C01--C05 podano lokalnie względem T0; C06--C08 w układzie SITE. Trzeci wymiar punktu kamery oznacza wysokość montażu, a punktu celu wysokość płaszczyzny obserwacji. Azymut liczony jest od osi wschodniej przeciwnie do ruchu wskazówek zegara. Pochylenia, odległości, granice kadru i odcinki obserwacji osoby są również zapisane w rejestrze i modelu.

\subsection*{Model optyczny i kontrola punktów}
Przyjęto projekcję prostoliniową otworkową 1920 $\times$ 1080, kwadratowe piksele i obrót obrazu 0$^\circ$. Korekcja dystorsji beczkowej jest włączona. W modelu obliczeniowym resztkowa dystorsja wynosi zero; po każdej stronie pozostawiono rezerwę 2\% całego wymiaru obrazu, czyli dopuszcza się współrzędne znormalizowane $|x|,|y|\leq0{,}96$. Jest to założenie projektowe, nie zmierzona kalibracja obiektywu. Karta producenta potwierdza zakresy HFOV 36--100$^\circ$, VFOV 20--53$^\circ$ i dostępność korekcji, a wybrane pary kątów wynikają z przyjętej projekcji. Wszystkie mieszczą się w obu zakresach.\cite{S-CAM}

Dla poziomego pola $\alpha$, pionowego $\beta$, wektorów osi kamery $\mathbf f$ (przód), $\mathbf r$ (prawo), $\mathbf u$ (góra), punktu $P$ i kamery $C$:
\[
\beta=2\arctan\left(\tan\frac{\alpha}{2}\frac{1080}{1920}\right),\qquad
D=\mathbf f\cdot(P-C)>0,
\]
\[
x=\frac{\mathbf r\cdot(P-C)}{D\tan(\alpha/2)},\qquad
y=\frac{\mathbf u\cdot(P-C)}{D\tan(\beta/2)}.
\]
Sektor jest prostokątną piramidą w 3D, dodatkowo ograniczoną poziomą odległością promieniową near/far oraz ścianami. Nie jest płaskim wycinkiem koła. Osie i cztery promienie narożne w SITE podaje eksport rejestru. B otrzymuje także parametry poniżej; nie ustala samodzielnie innego VFOV ani pochylenia.

\OpticsTable

Szerokość kadru w tabeli kamer to $w=2D\tan(\alpha/2)$ na płaszczyźnie prostopadłej do osi przez punkt celu; $1920/w$ jest gęstością w tej płaszczyźnie. Dla rzeczywistych otworów i ujęć osoby obliczono osobno rzut fizycznego odcinka poprzecznego 0,6 m, uwzględniając skośną perspektywę. Sprawdzane są oba krańce wysokości i odcinków. Własnym celem D01/D02 i PATH1 jest 250 px/m; dla przeglądu D03/D04 przyjęto 80 px/m, a G01 100 px/m. Przegląd służy powiązaniu zdarzenia z obrazem, bez obietnicy identyfikacji twarzy.

\OpticalVerificationTable

D01 obejmuje wyłącznie C01: zawężono HFOV do 40$^\circ$, pochylenie do 8$^\circ$, a granicę daleką ustawiono na 10,5 m. Kontrola dotyczy całego światła drzwi 1,2 m na wysokości 1,2--2,0 m, w odległości poziomej 9,5 m. C02 jest ujęciem tyłu na PATH1 i nie jest kamerą weryfikacji D01. D02 obejmuje C04. C05 ma skorygowany kierunek 49$^\circ$, HFOV 72$^\circ$ i granicę 13 m, aby objąć pełne światła D03/D04. G01 jest widoczna z C08 w płaszczyźnie otwartej bramy; zamknięta siatka może zasłaniać osobę po drugiej stronie. Weryfikację karty i użytkownika bramy wykonuje patrol.

C01--C04 obserwują prosty odcinek PATH1 w R102 z przodu, z tyłu i z obu boków. Kadry kończą się na ścianach pomieszczenia; C05 nie zapewnia obserwacji przez zamknięte D03/D04. Parking C07 daje przegląd zdarzeń, z możliwym zasłonięciem przez pojazdy. C08 obejmuje wydzielony pas dla pojazdów, a nie identyfikuje tablic. Wyłączono mikrofony, obserwację WC, stanowisk pracy i sąsiednich działek.\cite{S-AOA,L-GDPR}

\subsection*{Zapis i pojemność}
Przyjęto H.265, 25 kl./s, nominalnie 4 Mb/s i ograniczenie maksymalne do 6 Mb/s na kamerę. Bilans retencji wykorzystuje wartość maksymalną, łącznie 48 Mb/s:
\[
S=\frac{8\cdot6\cdot10^6}{8}\cdot7\cdot86400
=\StorageRaw{}\ \mathrm{TB}.
\]
Z rezerwą 25\% potrzeba \StorageNeed{} TB. Cztery dyski HAT3300-8T po 8 TB w Synology DS923+, RAID 6, dają po odłożeniu 10\% na narzuty i wolne miejsce \StorageHave{} TB. Nagrania zapisuje AXIS Camera Station Pro na udziale SMB3; NAS nie jest drugim VMS. RAID toleruje awarię dysków, lecz nie jest kopią materiału dowodowego. Jednostki są dziesiętne. Wymagana jest próba odtwarzania podczas odbudowy macierzy.\cite{S-NAS,S-VMS}

Każda kamera ma kartę 128 GB. Po odłożeniu 20\% przy 6 Mb/s bufor wynosi \SDHours{} h wobec celu 24 h. Odtworzenie luki zapisowej trzeba sprawdzić w funkcji failover recording VMS. Karta pomaga przy przerwie transmisji z zachowanym zasilaniem. Awaria przełącznika lub wyłączenie UPS usuwa także PoE i zatrzymuje zapis lokalny.
